\def\ChapterCode{INK}
\chapter
[\CO{[INK]} Infektions- und Entzündungskrankheiten]
{Infektions- und Entzündungskrankheiten}
\label{chp:INK}

\ChapterIntro
{%
~~~
}
{%
\begin{description}
  \item[Maintainer:] Sebastian Gabriel
  \item[Autoren:] Diverse
  \item[Reviewer:] Standard-Reviewprozess
  \item[Version:] {\DocVersion} ({\DocVersionInfo})
  \item[SHA1:] {\DocGitCommit}
\end{description}

{\TxTeilkapitelAASS}
}

\clearpage
%%%%%%%%%%%%%%%%%%%%%%%%%%%%%%%%%%%%%%%%%%%%%%%%%%%%%%%%%%%
%%% Infektionskrankheiten                               %%%
%%%%%%%%%%%%%%%%%%%%%%%%%%%%%%%%%%%%%%%%%%%%%%%%%%%%%%%%%%%

\section{Alltägliche Erkrankungen}


%%%%%%%%%%%%%%%%%%%%%%%%%%%%%%%%%%%%%%%%%%%%%%%%%%%%%%%%%%%
%%%%%%%%%%%%%%%%%%%%%%%%%%%%%%%%%%%%%%%%%%%%%%%%%%%%%%%%%%%
%%%%%%%%%%%%%%%%%%%%%%%%%%%%%%%%%%%%%%%%%%%%%%%%%%%%%%%%%%%
\subsection{Grippe}
\index{Grippe}
\index{Influenza}
\index{Influenza}
\label{topic:grippe}
\label{topic:influenza}

\HeadingSub{\Lang{Term.} \T{Influenza}}

\Definition{Eine schwere, durch Viren verursachte
Infektionskrankheit.}

\TextSlide{\TxBeschreibung}
{%
Die Grippe ist eine \E{virale} Erkrankung, welche je nach
Patient und Virenstamm (s.u.) für den Körper sehr belastend
sein kann und einen sehr schweren Verlauf nehmen und auch
tödlich enden kann. Gefürchtet sind \EE{Komplikationen} wie
Lungen- oder Herzmuskelentzündungen oder weitere
(Folge-)Infektionen. Besonders gefährdet sind ältere und
immunsupprimierte Personen.

\Fazit{Da es sich bei der Grippe um eine durch Viren
verursachte Krankheit handelt sind \EE{Antibiotika
wirkungslos}.}

Die Grippe tritt besonders in der kalten Jahreszeit,
meistens als \index{Epidemie}\EE{Epidemie} auf. Es kann auch
zum Ausbruch einer weltweiten \index{Pandemie}Pandemie
kommen, \Z{z.\,B. Spanische Grippe (1918, 25~Mio~Tote),
Asiatische Grippe (1957, 1~Mio~Tote) und  Hongkong-Grippe
(1968,
800.000~Tote)\footnote{http://de.wikipedia.org/wiki/Influenza}}.
\marginlabel{\cite{wikiSpanischeGrippe}}

In Österreich erkranken pro Jahr ca. 380.000 Personen an der
Influenza, ca. 4.500 Personen müssen stationär behandelt werden und
ca. 120 Personen sterben an den
Folgen\cite{Strauss0000}.

%\footnote{http://www.bmg.gv.at/cms/site/standard.html?channel=CH0742\&doc=CMS1075899626901}.

\Fazit{Die Grippe tritt in Wellen auf und kann tödlich
enden.}
}{
\begin{itemize}
  \item Viral
  \item Sehr belastend für Körper
  \item Schwere Komplikationen möglich
  \item Epidemien (\enquote{Grippewelle})
\end{itemize}
}{}{}{}




\Layout{60}{\TxSymptome}
{%
Typisch für die Grippe ist das \EE{rasche, hohe Anfiebern
(\VonBis{38}{40}\DegC*) mit Schüttelfrost, 
Kopf-, Hals-, Muskel- und Gelenksschmerzen, Husten und stark
reduziertem Allgemeinzustand}, sowie lang dauernde
Abgeschlagenheit nach Ende der Erkrankung. Bei schon vorher
geschwächten Patienten kann der Verlauf weniger spektakulär
aussehen (z.\,B. kaum Fieber), das Risiko einer
Komplikation ist allerdings natürlich höher.
}
{
\begin{itemize}
  \item Hohes Fieber
  \item Stark reduzierter Allgemeinzustand
  \item Gelenksschmerzen
\end{itemize}
}
{./images/gabriel-sebastian-aass/fieberthermometer.jpg}
{}
{GABS.}


\Text{Stämme}
{
Grippeviren kommen in verschiedenen
Stämmen (\Code{A}, \Code{B} und \Code{C}) vor, welche
wiederum durch unterschiedliche Arten von
\Z{Hämagglutinin} (\Code{H}) und \Z{Neuramidase}
(\Code{N}) unterschieden werden. Dadurch kommen so
klingende Namen wie \enquote{\Code{A(H5N1)}} oder
\enquote{\Code{A(H1N1)}} zu Stande. Je nach Unterart verhalten
sich die Viren anders und man beobachtet Unterschiede
hinsichtlich der Ansteckungsgefahr, Schwere der Erkrankung
usw.
\Z{}
}
{

}
{}
{}
{}

\TextSlide{Schutzimpfung}
{%
Gegen die Grippe wird eine \EE{Schutzimpfung} angeboten. Sie
basiert auf der \E{Voraussage}, welcher Stamm in der jeweils
bevorstehenden Grippewelle vorherrschen wird. Sie wirkt
daher nur für ein Jahr (eine Saison) und muss jede Saison
wiederholt werden. Die Impfung bietet keinen
hundertprozentigen Schutz, wird aber trotzdem für
Risikopatienten und Personen die einem erhöhten
Infektionsrisiko ausgesetzt sind (darunter fällt auch
Personal im Gesundheitsbereich) empfohlen. Die Impfung hat
einen schlechten Ruf, da Geimpfte häufig über leichtes
Fieber oder leichte Verkühlungen ein paar Tage nach der
Impfung klagen.
}{
\begin{itemize}
  \item Schutzimpfung für \E{vorhergesagten} Virenstamm der Saison
  \item Kein 100\,\%iger Schutz
  \item Vorhersage kann falsch sein
\end{itemize}
}
{}
{}
{}

\cite{Gallaher2009,H1N1InvTeam2009,Peiris2009,Wolf2009,Cutler2009}



\TextSlide{Schweinegrippe, Vogelgrippe}
{\Z{%
Grippeviren befallen nicht nur den Menschen, auch
(Wasser-)Vögel und andere Tiere erkranken an der Grippe.
Normalerweise erfolgt keine Übertragung vom Tier zum
Menschen. Überschreitet ein Erreger diese Grenze wird es
gefährlich: Die Tierwelt beinhaltet dann einen nahezu
unbegrenzten Vorrat an Keimen (Keimreservoir). Dieser kann 
praktisch nicht kontrolliert werden, schließlich lassen
sich Zugvögel nicht so leicht unter Quarantäne stellen
\ldots 

Glücklicherweise entwickelten sich in letzter Zeit nur
Viren, welche zwar vom Tier auf den Menschen und dann
weiter auf den Menschen übertragbar waren, jedoch waren sie
(bis jetzt) nicht übermäßig infektiös oder zeigten keinen
besonders schweren Verlauf. 

\subparagraph{Die Gefahr} Die
Tier\,\Sign{→}\,Mensch\,\Sign{→}\,Mensch-Übertragung stellt eine ständige
Gefahr dar. 

Ein weiteres Problem ist jedoch die mögliche Kreuzung mit
Stämmen die sehr infektiös und schwer verlaufend sind. Dies
kann passieren wenn ein Patient mit mehreren Stämmen
gleichzeitig infiziert ist. Daher ist es wichtig die
Infektionen so rasch als möglich einzudämmen.
}

\Fazit{Die große Gefahr ist die Entwicklung eines hoch
infektiösen und schwer verlaufenden Stammes, welcher eine
Tier-Mensch-Mensch-Übertragung aufweist.}
\Fazit{Eine rasche Isolierung der Infizierten soll die
Wahrscheinlichkeit einer bösartigen Kreuzung minimieren!}

\subparagraph{Neue Grippe (\enquote{Schweinegrippe})} \Z{Die jüngste
Entwicklung war die Übertragung des Stammes \Code{A(H1N1)}
vom Schwein auf den Menschen in Mexiko im Frühjahr 2009. Der
neue Stamm breitete sich sehr rasch um die ganze Welt aus,
\E{die Erkrankung verlief jedoch recht mild}. Bemerkenswert
war jedoch, dass -- anders als bei einer \enquote{normalen
Grippe} üblich -- besonders junge Menschen infiziert wurden. 

Man darf gespannt sein was die Zukunft bringt \ldots}
}
{%
\Z{%
\begin{itemize}
  \item Keimreservoir: Tiere
  \item Gefahr: Kreuzung von Virenstämmen
  \item Gefahr: Mensch-Mensch-Übertragung
\end{itemize}
}
}
{}
{}
{}


\Text{Wichtige Differentialdiagnosen}{%
Wie bei jedem fieberhaften Infekt muss an die Meningitis
gedacht werden und auf Meningismus-Zeichen geprüft werden
(\ref{topic:meningismuszeichen}).}
{
\begin{itemize}
  \item Meningitis \Sign{→} Meningismus-Zeichen
\end{itemize}
}{}{}{}


% \begin{tabulary}{\linewidth}{lLL}\toprule
% 
% \TabHeading{Übertragung}
% &
% \TabHeading{~}
% &
% \TabHeading{Einschätzung}
% \\\midrule
% Mensch \Sign{→} Mensch
% &
% &
% \\
% Tier \Sign{→} Mensch
% &
% &
% \\
% Tier \Sign{→} Mensch \Sign{→} Mensch
% &
% &
% \\
% Tier \Sign{→} Mensch \Sign{→} Mensch
% &
% hohe Ansteckungsrate
% 
% schwerer Verlauf
% &
% sehr schlimm
% \\\bottomrule
% \end{tabulary}

%%%%%%%%%%%%%%%%%%%%%%%%%%%%%%%%%%%%%%%%%%%%%%%%%%%%%%%%%%%
%%%%%%%%%%%%%%%%%%%%%%%%%%%%%%%%%%%%%%%%%%%%%%%%%%%%%%%%%%%
%%%%%%%%%%%%%%%%%%%%%%%%%%%%%%%%%%%%%%%%%%%%%%%%%%%%%%%%%%%
\subsection{Grippaler Infekt}
\index{grippaler Infekt}
\index{Verkühlung}
\index{Erkältung}


\HeadingSub{\Lang{Syn.} Verkühlung, Erkältung} 

\Text{\TxBeschreibung}
{%
Die Verkühlung ist eine Infektion der oberen Luftwege mit
Viren, allerdings mit anderen als bei der \enquote{echten Grippe}
(Influenza). Die Symptome ähneln denen der Grippe,
allerdings sind sie bei weitem nicht so stark ausgeprägt. Im
Vordergrund stehen Halsschmerzen, Husten, rinnende oder
verstopfte Nase und eine eher leicht erhöhte Temperatur.
}{}{}{}{}

\Text{Wichtige Differentialdiagnosen}{% Wie bei jedem
fieberhaften Infekt muss an die Meningitis gedacht werden
und auf Meningismus-Zeichen geprüft werden
(\ref{topic:meningismuszeichen}).}
{
\begin{itemize}
  \item Meningitis \Sign{→} Meningismus-Zeichen
\end{itemize}
}{}{}{}


%%%%%%%%%%%%%%%%%%%%%%%%%%%%%%%%%%%%%%%%%%%%%%%%%%%%%%%%%%%
%%%%%%%%%%%%%%%%%%%%%%%%%%%%%%%%%%%%%%%%%%%%%%%%%%%%%%%%%%%
%%%%%%%%%%%%%%%%%%%%%%%%%%%%%%%%%%%%%%%%%%%%%%%%%%%%%%%%%%%
\subsection{Lungenentzündung}
\label{topic:pneumonie}
\label{topic:lungenentzuendung}

\HeadingSub{\Lang{Term.} \T{Pneumonie}}

Streng genommen ist die Lungenentzündung eine Erkrankung der
Atemwege. Meistens steht bei den Symptomen eher die
entzündliche Komponente im Vordergrund (Fieber, allgemeine
Schwäche, \ldots).

\TextSlide{Ursachen}
{%
Eine Lungenentzündung kann grundsätzlich durch eine
Infektion durch Bakterien, Viren oder
Pilze\ZFootnote{(häufig bei HIV/Immunsuppression)} oder eine
chemische Schädigung, z.\,B. als Folge einer Aspiration von
Magensaft
  \index{Pneumonie!Aspiration} entstehen.

Höheres Lebensalter begünstigt Infektionen wie z.\,B.
Pneumonien, da ältere Patienten ein schwächeres Abwehrsystem
haben.
}{
\begin{itemize}
  \item Infektiös: Bakterien, Viren, Pilze
  \item Chemische Schädigung: z.\,B. nach
  \index{Pneumonie!Aspiration}Aspiration
\end{itemize}
}{}{}{}


\TextSlide{\TxSymptome}
{%
Die Symptome sind oft nicht charakteristisch, eine Diagnose
kann meist erst im Krankenhaus gestellt werden. Der Patient
hat oft einen \E{reduzierten Allgemeinzustand} und hustet,
manchmal klagt er über eine leichte Atemnot. \E{Fieber}
kann, muss aber nicht vorhanden sein. Besonders ältere
Menschen fiebern selbst bei schwereren Erkrankungen kaum an.

\bigskip

\bigskip
}{
\begin{itemize}
  \item uncharakteristisch:
  \begin{itemize}
    \item leichte Atemnot, Husten,  $\pm$ Fieber
  \end{itemize}
  \item Reduzierter Allgemeinzustand (AZ)
\end{itemize}
}{}{}{}



% M %%%%%%%%%%%%%%%%%%%%%%%%%%%%%%%%%%%%%%%%%%%%%%%%%%% M %
% Pneumonie
\ProcedurePrintDefault{IJ18090C}

\clearpage % TEMP temp clearpage
\section{Nicht alltägliche Erkrankungen}

%%%%%%%%%%%%%%%%%%%%%%%%%%%%%%%%%%%%%%%%%%%%%%%%%%%%%%%%%%%
%%%%%%%%%%%%%%%%%%%%%%%%%%%%%%%%%%%%%%%%%%%%%%%%%%%%%%%%%%%
%%%%%%%%%%%%%%%%%%%%%%%%%%%%%%%%%%%%%%%%%%%%%%%%%%%%%%%%%%%


\subsection{Tuberkulose}
\label{topic:tuberkulose}

\HeadingSub{\Z{Mykobacterium Tuberculosis, \enquote{Schwindsucht}.}
\Lang{Abk.} \enquote{\T{Tbc}}}

\TextSlide{Beschreibung}
{
Die Tuberkulose kann praktisch alle Organe befallen, der
Verlauf wird vor allem davon bestimmt, welches
Organ befallen ist \Z{(Lunge, Skelett, Darm, Haut,
Genitalien, ...)}. Ausgangspunkt ist fast immer die \EE{Lunge}.

Die Übertragung erfolgt vor allem \EE{aerogen}, aber unter
Umständen auch oral, durch Hautkontakt oder plazentär. 

Zunehmend treten auf Antibiotika resistente Erreger auf \Z{(MDR-TB,
XDR-TB)\footnote{Quelle: AGES}}.
% TODO LIT, HYG: MDR/XDR-TB Quelle bei AGES


\TODO{GABS}{2010-04-18}{Bild: Bild Lungenröntgen TB-Patient als Eye-Catcher.
--> Hauer}{}
    
\Z{Die Inkubationszeit beträgt 4-12 Wochen.}
}{}
% {./images/AASdev20090228-img98.jpg}
{./icons/under-construction.png}
{\AuthNotesPicLegalNotOk{img98}{Durch Foto von Hauer
ersetzen}Röntgenbild eines TB-Pat\-ient\-en.}{}


\TextSlide{\TxSymptome}
{%
Es dominieren vor allem Allgemeinsymtptome wie erhöhte
Temperatur, Abgeschlagenheit, Nachtschweiß,
Lymphknotenschwellung und Husten. Auffallend ist manchmal
die Schwere des Hustens, Blutbeimengungen gelten als
Alarmsignal.

In der Anamnese ist mitunter eine bereits überstandene
Tuberkulose erkennbar.
}{
\begin{itemize}
  \item Husten, Bluthusten
  \item Erhöhte Temperatur
  \item \EE{Nachtschweiß, Lymphknotenschwellung}
  \item Anamnese!
\end{itemize}
}
{}
{}
{}

\SlideCompensation{1}

\TextSlide{Offene Tuberkulose}
{%
\index{Offene Tb}%
Wenn der Infektionsherd in der Lunge \E{Luftkontakt} hat,
spricht man von der \T{Offenen Tuberkulose}. Der Patient ist
dann infektiös, eine Tröpfcheninfektion ist möglich. Die
Verwendung von Schutzmasken ist anzuraten.

Es ist äußerlich \E{nicht sicher erkennbar} ob der Patient
an einer offenen oder geschlossenen Tuberkulose leidet.
Bluthusten wäre allerdings ein deutlich Hinweiszeichen.
}
{
\begin{itemize}
  \item Infektionsherd hat Luftkontakt
  \item Tröpfcheninfektion möglich
\end{itemize}
}
{}
{}
{}


\Z{
\paragraph{Ansteckungsgefahr}
Bei Einhaltung der Schutzmaßnahmen (s. Spezielle Maßnahmen)
ist eine Ansteckung von sonst gesunden und immunkompetenten
Personal unwahrscheinlich. Für eine Ansteckung  wäre ein ca.
8-stündiger Kontakt erforderlich.\footnote{National
Institute for Clinical Excellence (NICE) Tuberculosis.
National clinical Guideline for Diagnosis, management,
prevention and control (2006)}
% TODO Fußnote -> Literatur
}



% M %%%%%%%%%%%%%%%%%%%%%%%%%%%%%%%%%%%%%%%%%%%%%%%%%%% M %
% Tuberkulose
\ProcedurePrintDefault{IA16091C}


\clearpage % TEMP temp clearpage
%%%%%%%%%%%%%%%%%%%%%%%%%%%%%%%%%%%%%%%%%%%%%%%%%%%%%%%%%%%
%%%%%%%%%%%%%%%%%%%%%%%%%%%%%%%%%%%%%%%%%%%%%%%%%%%%%%%%%%%
%%%%%%%%%%%%%%%%%%%%%%%%%%%%%%%%%%%%%%%%%%%%%%%%%%%%%%%%%%%
\subsection{HIV und AIDS}
\index{HIV}%
\index{Humanes Immundefizienz Virus}%
\label{topic:hiv}%
\label{topic:aids}%

\HeadingSub{%
\T{HIV}: \I{Humanes Immundefizienz Virus}; 
\T{AIDS}: \I{Aquired Immune Deficiency Syndrom},
\E{erworbenes Immunschwäche-Syndrom}}
\bigskip
\Layout{60}{HIV und AIDS}
{%
Das \T{HI-Virus} befällt vor allem die
Zellen des Abwehrsystems. Es vermehrt sich in ihnen, setzt sie
außer Funktion und zerstört sie schließlich. Das
körpereigene Abwehrsystem kann HIV nicht aus dem Körper
    entfernen, es kommt zu einer \II{Immunschwäche}.

HIV verursacht das \T{AIDS} (\I{Aquired
Immune Deficiency Syndrom}, erworbenes
Immunschwäche-Syndrom), welches die eigentliche Erkrankung
darstellt. Dabei kommt es durch das \E{geschwächte
Immunsystem} zu \EE{Folgeerkrankungen}, an denen der Patient
schwer erkranken und versterben kann:
\begin{itemize}
  \item \E{Infektionen} mit normalerweise harmlosen
  Erregern
  \begin{itemize}
    \item \Z{Pilze, \enquote*{fakultative} Erreger,
    \enquote*{Opportunisten}, seltene Arten der
    Lungenentzündung}
  \end{itemize}
  \item \E{Tumore}
\end{itemize}
Abhängig von der durchgeführten Therapie bricht das AIDS
einige Jahre nach der eigentlichen HIV-Infektion aus.
}{
\begin{itemize}
    \item HIV
    \item AIDS $\rightarrow$ Folgeerkrankungen
    \begin{itemize}
      \item Infektionen
      \item Tumore
    \end{itemize}
\end{itemize}
}
{./images/hygiene/Niaid-hiv-virion-mod.jpg}
{Schema eines HI-Virus.}
{US National Institute of Health, PD-USGov-HHS-NIH.}









\TextSlide{Übertragung}
{%
Übertragen wird das HI-Virus v.\,a. durch Körperflüssigkeiten
wie Blut, Blutprodukte, Sperma und Scheidensekret, Liquor oder Muttermilch.
\E{\enquote{Ungeschützter Sex}}, Transfusionen und
\E{Nadelstichverletzungen} gelten als klassische Risiken.
Es erfolgt \EE{keine Übertragung durch normale
Sozialkontakte} wie Händeschütteln, wohnen im
gemeinsamen Haushalt, gemeinsame Benutzung von Geschirr und
Toilettenanlagen, etc.
}{
\begin{itemize}
  \item Körperflüssigkeiten
  \item \E{Nicht} durch normale Sozialkontakte
\end{itemize}
}
{}
{}
{}
\SlideCompensation{1}

\TextSlide{Test}
{%
Erst nach \VonBis{6}{12} Wochen sind Antikörper nachweisbar.
Das bedeutet, dass der \EE{HIV-Test nicht \enquote{aktuell},
sondern \enquote{12 Wochen blind}} ist!
}
{%
\begin{itemize}
  \item Erst einige Wochen nach Infektion aussagekräftig
\end{itemize}
}
{}
{}
{}
\SlideCompensation{1}

\TextSlide{Therapie}
{%
Derzeit ist \EE{keine Heilung} möglich, jedoch
  kann eine \E{deutliche Lebensverlängerung} durch
  Medikamente erreicht werden, indem der Krankheitsausbruch
  (AIDS) verzögert wird.

}
{
\begin{itemize}
  \item Keine Heilung, aber Verzögerung des
  Krankheitsausbruchs möglich
\end{itemize}
}
{}
{}
{}
\SlideCompensation{}

\Text{Prophylaxe}
{%
Allgemeine Selbstschutzmaßnahmen bei der
Patientenversorgung sind ausreichend. Dazu gehört das
Tragen von Untersuchungshandschuhen und die Vermeidung
des Kontakts mit Körperflüssigkeiten. Kommt es zu einer
Nadelstichverletzung kann innerhalb von zwei Stunden eine
\I{Postexpositionsprohylaxe} (\I{PEP}) verschrieben werden.
Bei Sexualkontakten stellt die Verwendung von Kondomen einen
effektiven, wenn auch keinen 100\PC*-igen Schutz dar.

}
{
\begin{itemize}
  \item Allgemeine Selbstschutzmaßnahmen 
  \item Postexpositionsprohylaxe bei Nadelstichverletzungen
  \item Kondome
\end{itemize}
}
{}
{}
{}



\clearpage
%%%%%%%%%%%%%%%%%%%%%%%%%%%%%%%%%%%%%%%%%%%%%%%%%%%%%%%%%%%
%%%%%%%%%%%%%%%%%%%%%%%%%%%%%%%%%%%%%%%%%%%%%%%%%%%%%%%%%%%
%%%%%%%%%%%%%%%%%%%%%%%%%%%%%%%%%%%%%%%%%%%%%%%%%%%%%%%%%%%
\subsection{Hepatitis}
\label{topic:hepatitis-virale}%
\label{topic:hepatitis-toxisch}%
\label{topic:hepatitis}%
\index{Hepatitis!infektiöse}%
\index{Hepatitis!toxische}%

\Definition{Leberentzündung durch verschiedene (virale)
Erreger (A -- E) oder durch sonstige
(chronische) Schädigung.}



\TextSlide{\TxBeschreibung{} und Ursachen}
{%%% Gerhard Gabriel <gerhard.gabriel@grissu.org> 
Hepatitis ist eine
entzündliche Erkrankung der Leber. Diese Entzündung kann
durch Erreger wie Viren, oder Mikroorganismen hervorgerufen
werden (infektiöse, virale Hepatitis). Aber auch regelmäßige
Schädigungen der Leber durch Giftstoffe wie Alkohol,
Medikamente \Z{(z.\,B. Paracetamol)}, etc. können zu einer
(chronischen) Entzündung führen. Auch ein länger bestehender
Rückstau der Gallenflüssigkeit, z.\,B. bei einer Verstopfung
des Gallenganges, kann zu einer Entzündung führen.

Je nach Ursache der Hepatitis kann es vorkommen, dass es zu
keiner Abheilung kommt und die Hepatitis dauerhaft bestehen
bleibt. Man spricht dann von der \T{chronischen
Hepatitis}.} 
{
\begin{itemize}
  \item Virusinfektionen (Hep. A, B, C, \ldots )
  \item Gallen-Abflussstau
  \item Vergiftungen Lebensmittel(Pilze)~/
  Medikamente (Paracetamol)~/ chron. Alkoholismus
  \item Chronische Verlaufsform möglich
\end{itemize}
}{}{}{}

\TextSlide{\TxSymptome}
{%
Die Symptomatik einer Hepatitis ist uncharakteristisch und
eher allgemein: Unspezifisches mehr oder minder ausgeprägtes
Krankheitsgefühl, ähnlich einem grippalen Infekt,
Appetitlosigkeit. Auch erhöhte Körpertemperatur und
Schmerzen im rechten Oberbauch kommen bisweilen vor. Die
Leber ist oft
vergrößert.
Im
fortgeschrittenen Stadium kann es zu einer \EE{Gelbfärbung}
des Augenweiß und einer gelblichen Verfärbung der Haut
(\I{Ikterus}) kommen. 

Bei einem chronischen
Verlauf kommt es zu einem Umbau von
funktionierendem Lebergewebe zu
funktionslosem Bindegewebe. Man spricht
dann von der
\T{Leberzirrhose}\label{topic:leberzirrhose}\index{Leberzirrhose}.
Schließlich kann es zu einem Leberversagen und in kurzer
Folge zum Tod kommen, da die Leber ein überlebenswichtiges
Organ ist. 
}{
\begin{itemize}
  \item uncharakteristische Allgemeinsymptome (grippale Symptome,
  Appetitlosigkeit)
  \item Schmerzen im re. Oberbauch
  \item Gelbfärbung Augen/Haut (Ikterus )
  \item Vergrößerte, schmerzhafte Leber
  \item Ikterus (Gelbsucht)
  \item Unterscheide \E{akut} und
  \E{chronisch}
  \begin{itemize}
    \item Chronisch: wenn Hepatitis nicht abheilt
  \end{itemize}
  \item Ende: Leberversagen, Tod
\end{itemize}
}{}{}{}



\paragraph{Verschiedene Arten der viralen Hepatitis}~

Abhängig vom jeweiligen Virus unterscheidet man
verschiedene Arten der viralen Hepatitis:

\begin{table}[H]
\FontTableBase
\begin{WidePar}
\Captionof{table}{Übersicht der wichtigsten viralen Hepatitis-Arten.}

\begin{tabular}{p{3cm}p{\linewidth-4cm}}\addlinespace\toprule\addlinespace
\noindent\TabHeading{Typ}
&
\noindent\TabHeading{Bemerkungen}
\\\addlinespace\midrule\addlinespace
\TabSub{Hepatitis-A} 

Hepatitis-A-Virus (HAV)
&
\begin{itemize}
  \item \Z{Inkubationszeit 10 - 40 Tage}
  \item typische Reisekrankheit (Mittelmeerraum, Südamerika, Orient),
  meist epidemisch,
  \item Dauer einige Wochen,
  \item kein Übergang in chronischen Verlauf
\end{itemize}

\\\addlinespace
\TabSub{Hepatitis-B} 

Hepatitis-B-Virus (HBV)
&
\begin{itemize}
  \item \Z{Inkubationszeit 4 - 12 Wochen,}
  \item \Z{meist sporadisches Auftreten}
  \item Übertragung durch Speichel, Blut(produkte),
  Sperma und andere Körperflüssigkeiten,
  \item chronischer Verlauf in ca. 5 \%
\end{itemize}
\\\addlinespace
\TabSub{Hepatitis C}

Hepatitis-C-Virus (HCV)
&
\begin{itemize}
  \item \Z{Inkubationszeit 6 - 12 Wochen,}
  \item häufig nach Bluttransfusion oder i.v. Drogenabusus,
  \item \EE{oft chronischer Verlauf} (in 70 -- 80~\% der
  Fälle)
\end{itemize}
\\\bottomrule
\end{tabular}
\end{WidePar}
\end{table}



\TextSlide{Impfung}
{%
\index{Impfschutz!Hepatitis}%
\index{Hepatitis!Impfschutz}%
Für die Hepatitis A und B ist eine Impfung verfügbar (z.\,B.
Twinrix{\texttrademark}). Für die Hepatitis C ist derzeit
noch keine Impfung regulär am Markt. Gegen die toxischen
Formen der Hepatitis gibt es natürlich auch keine Impfung.

Achtung: Non-Responder möglich!
% FEATURE Hepatitis Laborwerte Interpretation einfügen.
}{
\begin{itemize}
  \item Für Hep. A \& B
  \item Nicht für Hep. C
\end{itemize}
}
{}
{}
{}


\Fazit{Für den Rettungsdienst sind besonders die Hepatitis
B und C wichtig: Hohe Infektiosität, chronischer Verlauf!}

\Fazit{Gegen Hepatitis C gibt es derzeit keine Impfung!}

\Fazit{Für Risikopersonal (dazu zählt auch Sanitätspersonal)
übernimmt i.\,d.\,R. die Unfallversicherungsanstalt die
Kosten für die Schutzimpfung.}




% \clearpage % TEMP temp clearpage
\section{Seltene, aber schwere Erkrankungen}

\subsection{Bakterielle Meningitis}
\index{Meningitis}
\label{topic:meningitis}

\Definition{Eitrige Entzündung der Hirnhäute aufgrund
einer Infektion mit Bakterien.}

\noindent Die bakterielle Meningitis ist eine gefürchtete
Erkrankung. Einerseits ist sie hoch ansteckend, andererseits
kann sie rasch tödlich enden. Gefährdet sind grundsätzlich
alle Altersgruppen, besonders aber Kinder, Jugendliche und
alte Menschen.
 
 
\TextSlide{\TxSymptome}
{%
\label{topic:meningismuszeichen}
\index{Meningismuszeichen!Nackensteifigkeit}
\index{Nackensteifigkeit}
Leitsymptome sind \EE{Fieber}, \EE{Nackensteifigkeit}
(Meningismus) und \EE{Lichtscheu}. Dies bezeichnet man als
\T{Meningismuszeichen}.

Darüber hinaus kann es zu anderen unspezifischen Symptomen
wie Kopfschmerzen, Übelkeit, Erbrechen, etc. kommen.

\E{Die Meningitis wird oft mit einem einfachen grippalen
Infekt oder einer Grippe verwechselt!}

\Z{Im fortgeschrittenem Stadium kann es zu massiven
Gerinnungsstörungen kommen, am Körper bilden sich dann
dunkle Hämatome. Einer der häufigsten Erreger sind die
Meningokokken, man spricht dann von einer
Meningokokken-Meningitis.}

}{
\begin{itemize}
  \item Fieber, Kopfschmerzen
  \item Nackensteifigkeit
  \item Lichtscheu
  \item Erbrechen, Übelkeit
  \item Hämatome am ganzen Körper (im Endstadium)
  \item Verwechslungsgefahr mit \enquote{normalem Infekt}!
\end{itemize}

}{}{}{}

\SlideCompensationNoParskip{1}

\opt{STATUS-EXPERIMENTAL}{

\Layout{22}{}
{
% Endstadium einer bakteriellen Meningitis.
% Die Flecken auf der Haut sind Blutungen in Folge einer nicht mehr
% funktionierenden Blutgerinnung. Der Tod ist bei dieser
% Patientin innerhalb weniger Stunden zu erwarten.

\TODO{GABS}{2010-04-18}{Bild: Meningitis von KAR + Beschreibung}{}
}{}
% {./images/meningitis-kind-waterhouse.png}
{./icons/under-construction.png}
{\label{fig:meningitis-baby}}
{unbekannt}
}





% M %%%%%%%%%%%%%%%%%%%%%%%%%%%%%%%%%%%%%%%%%%%%%%%%%%% M %
% Meningitis
\ProcedurePrintDefault{IG03091C}


\subsection{Wundstarrkrampf (Tetanus)}
\index{Tetanus}
\index{Wundstarrkrampf}
\label{topic:Wundstarrkrampf}

\TextSlide
{\TxBeschreibung}
{Der Erreger des Wundstarrkrampfes\ZFootnote{Sporen von
\I{Clostridium tetani}} kommt \E{überall in der Umgebung} vor. Bei
Verletzungen können durch Kontamination diese Sporen in die
Wunde kommen. Bei ungünstigen Wundverhältnissen
(\E{anaerobe} Verhältnisse) kann es dann zu einer
Infektion kommen. Der Erreger produziert
einen Giftstoff (Toxin), welches einen tonischen Krampf der
Skelettmuskulatur auslöst. \Z{Die Inkubationszeit beträgt
ca. \VonBis{4}{21} Tage. Erste Vorzeichen sind Müdigkeit und
Inappetenz. \cite{HahnMikrobio6}}

}
{
\begin{itemize}
  \item Erreger kommt überall vor
  \item Produziert Toxin $\rightarrow$ Muskellähmung
\end{itemize}
}
{}
{}
{}



\Layout{20}{\TxSymptome}
{%
Es kommt zu tonischen \EE{Krämpfen der
quergestreiften Muskulatur}, oft beginnend bei der
Kaumuskulatur \Z{(Kiefersperre, Teufelsgrinsen
(\I{Risus sardonicus}))}. Der Patient hat starke Schmerzen
und entwickelt in Folge der Krämpfe eine \EE{Atemlähmung}. 
Die Fortschritte der
Intensivmedizin haben die Behandlungsmöglichkeiten
verbessert, trotzdem sterben noch etwa \VonBis{20}{30}\PC*
der Tetanuserkrankten. Der Wert der
Impfung ist daher nach wie vor groß.
}
{
\begin{itemize}
  \item Krämpfe der Skelettmuskulatur
  \item Schmerzen, Atemlähmung
  \item Hohe Letalität
\end{itemize}
}
{./images/hygiene/6373_lores_tetanus.jpg}
{Patient mit Tetanus.}
{Public Health Image Library, ID\#: 6373, Public Domain.}

\TextSlide{Schutzimpfung}
{
\index{Impfschutz!Tetanus}
\index{Tetanus!Impfschutz}
\label{topic:tetanus-impfschutz}
Es existiert eine Schutzimpfung, diese muss dem Patienten
\E{rechtzeitig} nach der Verletzung gegeben werden. Der
Impfschutz beträgt normalerweise ca \E{10 Jahre}, bei einer
Verletzung wird sicherheitshalber allerdings schon bei einer
\EE{länger als 5 Jahre} zurückliegenden Impfung eine
neuerliche Gabe verabreicht.

Auch bei kleinen Verletzungen muss nach dem Impfschutz
gefragt werden, dies muss auch unbedingt \EE{dokumentiert
werden!} Ist das Datum der letzten Impfung nicht sicher
erinnerlich soll der Patient unverzüglich zu Hause im
Impfpass nachschauen. Erforderlichenfalls soll der Patient
unverzüglich eine Unfallabteilung aufsuchen.
Die entsprechende Aufklärung muss am Revers bestätigt werden!
}
{
\begin{itemize}
  \item Impfschutz \VonBis{5}{10} Jahre
  \item Bei Verletzung: Wenn letzte Impfung >\,5~Jahre:
  Auffrischung
\end{itemize}
}
{}
{}
{}

\Cave{Eine mangelhafte Aufklärung oder Dokumentation
bezüglich Tetanus und dessen Impfschutz ist grob fahrlässig und ein unverzeihbarer
Kunstfehler.}




%%%%%%%%%%%%%%%%%%%%%%%%%%%%%%%%%%%%%%%%%%%%%%%%%%%%%%%%%%%
%%%%%%%%%%%%%%%%%%%%%%%%%%%%%%%%%%%%%%%%%%%%%%%%%%%%%%%%%%%
%%%%%%%%%%%%%%%%%%%%%%%%%%%%%%%%%%%%%%%%%%%%%%%%%%%%%%%%%%%
\subsection{Sepsis}
\label{topic:sepsis}

Siehe auch: Septischer Schock und Toxischer Schock,
Kap.\,\ref{topic:schock-septischer}.

\begin{itemize}
  \item Wenn Keime und deren Produkt außer Kontrolle geraten,
  erfolgt Amoklauf des Körpers:
  \item Ungehemmte Freisetzung von Stoffen des Entzündungs-, Gerinnungs-
  und Immunsystems
  \item Auch ohne Infektion möglich:
  \begin{itemize}
    \item \Z{Systemic Inflammatory Response Syndrome (\Abk{SIRS})}
    \item nach schwerem Trauma, Verbrennungen, Gewebshypoxie
  \end{itemize}
\end{itemize}


\Text{\TxSymptome}
{%
Die Symptome sind sehr von Fortschritt und
körperlichem Abbau abhängig. Grundsätzlich:


\begin{itemize}
  \item Sehr stark reduzierter Allgemeinzustand
  \item Fieber ({\textgreater}\,38{\DegC}) oder
  Hypothermie ({\textless}\,36{\DegC})
  \item Herzfrequenz {\textgreater}\,90\nicefrac{}{\Min}
  \item Atemfrequenz {\textgreater}\,20\nicefrac{}{\Min} oder
  Hyperventilation
  \item unkontrollierbare Gerinnung/Blutung
  \item Hypotonie ${\rightarrow}$ Septischer Schock
  \item Organversagen
\end{itemize}

}{}{}{}{}


\ChapterEnd